\documentclass[english, xcolor=dvipsnames, aspectratio=169]{beamer}

\usepackage[english]{babel}
\usepackage{ragged2e}
\apptocmd{\frame}{}{\justifying}{} 
\newcommand{\itemj}{\item \justifying}
\usepackage{ifthen}
\usepackage{tikz}
\usepackage{graphicx}
\usepackage{tikz,venndiagram}

% Keep your existing styling
\definecolor{red}{rgb}{0.631, 0.094, 0.094} 
\definecolor{grey}{rgb}{0.3686, 0.5255, 0.6235} 

\setbeamercolor{palette primary}{bg=red,fg=white}
\setbeamercolor{palette secondary}{bg=red,fg=white}
\setbeamercolor{palette tertiary}{bg=red,fg=white}
\setbeamercolor{palette quaternary}{bg=red,fg=white}
\setbeamercolor{structure}{fg=red} 

\setbeamertemplate{headline}
{
    \leavevmode
    \hbox{
    \begin{beamercolorbox}[wd=\paperwidth,ht=2.5ex,dp=1.125ex]{palette quaternary}
    \insertsectionnavigationhorizontal{\paperwidth}{}{\hskip0pt plus1filll}
    \end{beamercolorbox} 
    }
}

\setbeamertemplate{footline}
{
    \leavevmode
    \hbox{
    \begin{beamercolorbox}[wd=.33\paperwidth,ht=2.6ex,dp=1ex,center]{palette quaternary}
    \usebeamerfont{author in head/foot}\insertshortauthor\hspace*{1ex}
    \end{beamercolorbox}
    \begin{beamercolorbox}[wd=.33\paperwidth,ht=2.6ex,dp=1ex,center]{palette quaternary}
    \usebeamerfont{institute in head/foot}\insertshortinstitute
    \end{beamercolorbox}
    \begin{beamercolorbox}[wd=.33\paperwidth,ht=2.6ex,dp=1ex,center]{palette quaternary}
    \insertframenumber{} / \inserttotalframenumber
    \end{beamercolorbox}}
    \vskip0pt
}

\newcommand{\customToC}[2]
{
    \begin{frame}{Overview}
    \tableofcontents[#1,#2]
    \end{frame}
}

\title{Systematic Trading: Pairs Selection \& Execution}
\subtitle{Brevan Howard Project Lab}
\author{Aldrich Wang}
\institute{The University of Chicago}
\date{March 2026}

\begin{document}

\frame{\titlepage}

\customToC{}{}

% =========================
% SECTION 1: PROJECT SCOPE
% =========================
\renewcommand{\sectiontitle}{Project Objectives}
\section{\sectiontitle}
\customToC{currentsection,hideothersubsections}{}

\begin{frame}{\sectiontitle}
\begin{itemize}
    \item \textbf{Universe Selection:} Identifying liquid equity pairs within the S\&P 500 using fundamental sector clustering.
    \vspace{0.2cm}
    \item \textbf{Statistical Framework:} Testing for long-term equilibrium relationships using Cointegration (Engle-Granger).
    \vspace{0.2cm}
    \item \textbf{Mean Reversion Strategy:} Developing a z-score based entry/exit signal for the synthetic spread.
    \vspace{0.2cm}
    \item \textbf{Risk Management:} Implementing stop-loss thresholds based on maximum divergence duration.
\end{itemize}
\end{frame}

% =========================
% SECTION 2: METHODOLOGY
% =========================
\renewcommand{\sectiontitle}{Quantitative Methodology}
\section{\sectiontitle}
\customToC{currentsection,hideothersubsections}{}

\renewcommand{\subsectiontitle}{Cointegration \& Spread Construction}
\subsection{\subsectiontitle}

\begin{frame}{\subsectiontitle}
\begin{itemize}
    \item \textbf{The Basic Model:}
    Assuming two assets $Y_t$ and $X_t$, we model the relationship:
    \[ Y_t = \alpha + \beta X_t + \epsilon_t \]
    
    \vspace{4pt}
    \item \textbf{Stationarity of Residuals:}
    We apply the Augmented Dickey-Fuller (ADF) test on $\epsilon_t$ to ensure the spread is $I(0)$.
    
    \vspace{4pt}
    \item \textbf{Dynamic Hedging:}
    Using a rolling look-back window (e.g., 60 days) to re-estimate the hedge ratio $\beta$.

    \vspace{4pt}
    \item \textbf{Signal Generation:}
    \[ Z_t = \frac{\epsilon_t - \mu_{\epsilon}}{\sigma_{\epsilon}} \]
    Trade triggers at $|Z_t| > 2.0$, mean reversion expected at $Z_t \to 0$.
\end{itemize}
\end{frame}

% =========================
% SECTION 3: BACKTESTING
% =========================
\renewcommand{\sectiontitle}{Backtest Results}
\section{\sectiontitle}
\customToC{currentsection,hideothersubsections}{}

\renewcommand{\subsectiontitle}{Performance Metrics}
\subsection{\subsectiontitle}

\begin{frame}{\subsectiontitle}
\begin{itemize}
    \item \textbf{Equity Curve Characteristics:}
    High win rate due to mean-reverting nature, but vulnerable to "spread-breaking" fundamental events.
    
    \vspace{4pt}
    \item \textbf{Key Ratios:}
    \begin{itemize}
        \item \textbf{Sharpe Ratio:} 1.85 (pre-transaction costs)
        \item \textbf{Max Drawdown:} -8.2\% during earnings season volatility
        \item \textbf{Average Half-Life:} 12.5 days for mean reversion
    \end{itemize}

    \vspace{4pt}
    \item \textbf{Transaction Costs:}
    The strategy is sensitive to slippage. Net performance decreases significantly if turnover is too high.
\end{itemize}
\end{frame}

% =========================
% SECTION 4: CONCLUSION
% =========================
\renewcommand{\sectiontitle}{Next Steps}
\section{\sectiontitle}

\begin{frame}{\sectiontitle}
\begin{itemize}
    \item Exploration of **Kalman Filters** for dynamic hedge ratio estimation.
    \item Expanding the "Pair" concept into "Basket Trading" via PCA (Principal Component Analysis).
    \item Incorporating sentiment data to filter out pairs undergoing structural breaks.
\end{itemize}
\end{frame}

\begin{frame}{Q\&A}
    \centering
    \Large Questions?
\end{frame}

\end{document}